\subsection{Motivation}

Automated currency trading is usually built from fixed rules. Technical indicators are computed from recent prices, a signal is derived from them by comparison or crossover, and a separate layer decides the size of the position \cite{chan_quantitative_2021, brandimarte_introduction_2018}. These systems are cheap to run and easy to inspect, which is why they are common. Their parameters are set once and stay fixed, so the system does not adapt when the market changes. That is the first of the two difficulties this work deals with.

Currencies were chosen over equities for practical reasons \cite{donnelly_art_2019}. The market trades continuously on weekdays, so a position is never carried through a daily close that cannot be traded. There are only seven major pairs, which makes it possible to test a design on all of them rather than on a sample selected by the author. Both directions are available without borrowing costs or short-selling restrictions, so a policy that sells is not penalised by mechanics unrelated to its signal. Every major pair also has the US dollar on one side, so the seven instruments are related without being interchangeable, and comparing results across them says something about the design rather than about one price series.

Reinforcement learning describes the same task without separating prediction from action. The agent observes the state of the market, chooses a position, and receives the profit or loss of that position as its reward \cite{sutton_reinforcement_2018}. It is never asked to forecast a price. It is asked to act, and the action is measured in money, which is how a trading system is judged in practice.

DDPG fits this task well. It is an actor-critic method for continuous action spaces \cite{lillicrap_continuous_2015}, so the output of the actor can be used directly as a signed position: the sign gives the direction and the magnitude gives the exposure. Much of the published work on trading uses value-based methods such as the Deep Q-Network \cite{mnih_human-level_2015, fischer_reinforcement_2018, hambly_recent_2023}, which need the action to be reduced to a small discrete set, normally buy, hold and sell. Position sizing is then decided outside the learned policy, by rules that were never trained against the reward.

The second difficulty is cost, and it is independent of the learning algorithm. Results in this field are often reported on bar data, with the spread fixed or ignored and the broker commission left out \cite{fischer_reinforcement_2018, pricope_deep_2021}. For a strategy that trades a few times per year this changes little. For a strategy that adjusts its position often it changes a lot, because the cost is paid on every adjustment while the profit per trade stays small. A system that is profitable without costs can lose money with them, and the difference does not appear unless the test charges what the account would pay \cite{lopez_de_prado_advances_2018}. The size of that difference is easy to underestimate. In this work, charging the broker commission on an otherwise identical run removed twenty percentage points of total return from a configuration that adjusted its position at every decision, and six percentage points once a threshold was added so that small adjustments were skipped. Both figures are reported in Section~
ef{Section:ExperimentalResults}.